\begin{lexlean}{Uorc.Core.RuntimeAdmissionBoundary}
\importmodule{Uorc.Core.RuntimeAdmission}
\importmodule{Uorc.Core.RuntimeLifetime}
\importmodule{Uorc.Core.RuntimeInstructionCheck}
\importmodule{Uorc.Core.RuntimeAdmissionOrder}
\importmodule{Uorc.Core.RuntimeProgram}
\importmodule{Uorc.Core.RuntimeResources}
\importmodule{Uorc.Core.Heap}
\importmodule{Uorc.Core.RuntimeValue}
\importmodule{Uorc.Core.RuntimeTable}
\importmodule{Uorc.Core.RuntimeByteBuilder}
\importmodule{Uorc.Core.RuntimeWork}
\title{Parameters}
\begin{semanticmodule}
inductive{
  name: Error,
  type_parameters: [],
  parameters: [],
  constructors: [
    {name: busy, fields: []},
    {name: wrongRole, fields: []},
    {name: staleSession, fields: []},
    {name: counterRegression, fields: []},
    {name: callerMismatch, fields: []},
    {name: foreignRelease, fields: []},
    {name: invalidTicketInventory, fields: []},
    {name: resource, fields: [Uorc.Core.RuntimeResources::Error]},
    {name: heap, fields: [Uorc.Core.Heap::HeapError]}
  ]
}

inductive{
  name: ProducerRole,
  type_parameters: [],
  parameters: [],
  constructors: [{name: wire, fields: []}, {name: search, fields: []}]
}

structure{
  name: AwaitingState,
  type_parameters: [],
  parameters: [],
  fields: [
    {name: budget, type: Uorc.Core.RuntimeResources::Budget},
    {name: program, type: Uorc.Core.RuntimeProgram::Program},
    {name: order, type: Uorc.Core.RuntimeAdmissionOrder::State},
    {name: pendingOrder, type: (Option Uorc.Core.RuntimeAdmissionOrder::State)},
    {name: pendingEvent, type: (Option Uorc.Core.RuntimeProgram::Event)},
    {name: instruction, type: Uorc.Core.RuntimeInstructionCheck::State},
    {name: parameterTypes, type: Uorc.Core.RuntimeTable::Table},
    {name: localTypes, type: Uorc.Core.RuntimeTable::Table},
    {name: instructions, type: Uorc.Core.RuntimeTable::Table},
    {name: lastUses, type: Uorc.Core.RuntimeTable::Table},
    {name: lifetimes, type: Uorc.Core.RuntimeTable::Table},
    {name: capturesA, type: Uorc.Core.RuntimeTable::Table},
    {name: capturesB, type: Uorc.Core.RuntimeTable::Table},
    {name: bucketSlots, type: Uorc.Core.RuntimeTable::Table},
    {name: initialRelease, type: Uorc.Core.RuntimeValue::Handle},
    {name: blockRecord, type: Uorc.Core.RuntimeValue::Handle},
    {name: instructionRecord, type: Uorc.Core.RuntimeValue::Handle},
    {name: bucketRecord, type: Uorc.Core.RuntimeValue::Handle},
    {name: literal, type: Uorc.Core.RuntimeValue::Handle},
    {name: auxFirst, type: Uorc.Core.RuntimeValue::Handle},
    {name: auxSecond, type: Uorc.Core.RuntimeValue::Handle},
    {name: control, type: Uorc.Core.RuntimeValue::Handle},
    {name: declaredResult, type: Uorc.Core.RuntimeValue::RuntimeType},
    {name: returnSlot, type: UInt64},
    {name: lifetimeCursor, type: (Option Uorc.Core.RuntimeLifetime::Cursor)},
    {name: pendingSlot, type: UInt64},
    {name: bucketIndex, type: UInt64},
    {name: bucketCount, type: UInt64},
    {name: bucketInitial, type: Bool},
    {name: task, type: Uorc.Core.RuntimeAdmission::Task},
    {name: failure, type: (Option Uorc.Core.RuntimeAdmission::Error)},
    {name: literalBuilder, type: Uorc.Core.RuntimeByteBuilder::Builder},
    {name: literalOffset, type: UInt64},
    {name: literalRemaining, type: UInt64},
    {name: role, type: ProducerRole},
    {name: lease, type: Uorc.Core.RuntimeValue::Handle},
    {name: expectedCaller, type: Uorc.Core.RuntimeProgram::CallerRoots},
    {name: minimumNextHandle, type: UInt64},
    {name: wasExhausted, type: Bool},
    {name: minimumWork, type: UInt64},
    {name: minimumOutput, type: UInt64},
    {name: expectedFrames, type: UInt64},
    {name: minimumPeak, type: UInt64}
  ]
}

structure{
  name: Taken,
  type_parameters: [],
  parameters: [],
  fields: [
    {name: state, type: AwaitingState},
    {name: arena, type: Uorc.Core.Heap::Arena},
    {name: usage, type: Uorc.Core.RuntimeResources::Usage},
    {name: caller, type: Uorc.Core.RuntimeProgram::CallerRoots}
  ]
}

structure{
  name: ReturnCursor,
  type_parameters: [],
  parameters: [],
  fields: [{name: taken, type: Taken}, {name: leaseLive, type: Bool}]
}

inductive{
  name: TakeEvent,
  type_parameters: [],
  parameters: [],
  constructors: [
    {name: taken, fields: [Taken]},
    {name: failed, fields: [Error, Uorc.Core.RuntimeAdmission::Cursor]},
    {name: retainedFailure, fields: [Error, Taken]}
  ]
}

inductive{
  name: ReturnEvent,
  type_parameters: [],
  parameters: [],
  constructors: [
    {name: pending, fields: [ReturnCursor]},
    {name: restored, fields: [Uorc.Core.RuntimeAdmission::Cursor]},
    {name: failed, fields: [Error, ReturnCursor]}
  ]
}

definition{
  name: sameRole,
  parameters: [first : ProducerRole, second : ProducerRole],
  result: Bool,
  body: (match var{name: first}
    | ProducerRole.wire => (match var{name: second}
      | ProducerRole.wire => true
      | ProducerRole.search => false)
    | ProducerRole.search => (match var{name: second}
      | ProducerRole.wire => false
      | ProducerRole.search => true)),
  axioms: [],
  executable: true
}

definition{
  name: sameNonProducerCaller,
  parameters: [
    role : ProducerRole,
    expected : Uorc.Core.RuntimeProgram::CallerRoots,
    actual : Uorc.Core.RuntimeProgram::CallerRoots
  ],
  result: Bool,
  body: (match role
    | ProducerRole.wire => ((((((true && (prim equal (project (project expected input) id) (project (project actual input) id) : Bool)) && (prim equal (project (project expected archive) id) (project (project actual archive) id) : Bool)) && (prim equal (project (project expected search) id) (project (project actual search) id) : Bool)) && (prim equal (project (project expected checkpoint) id) (project (project actual checkpoint) id) : Bool)) && (prim equal (project (project expected certificate) id) (project (project actual certificate) id) : Bool)) && (prim equal (project (project expected verification) id) (project (project actual verification) id) : Bool))
    | ProducerRole.search => ((((((true && (prim equal (project (project expected input) id) (project (project actual input) id) : Bool)) && (prim equal (project (project expected archive) id) (project (project actual archive) id) : Bool)) && (prim equal (project (project expected checkpoint) id) (project (project actual checkpoint) id) : Bool)) && (prim equal (project (project expected certificate) id) (project (project actual certificate) id) : Bool)) && (prim equal (project (project expected verification) id) (project (project actual verification) id) : Bool)) && (prim equal (project (project expected wire) id) (project (project actual wire) id) : Bool))),
  axioms: [],
  executable: true
}

definition{
  name: rootInventory,
  parameters: [state : AwaitingState, caller : Uorc.Core.RuntimeProgram::CallerRoots],
  result: (List Uorc.Core.RuntimeValue::Handle),
  body: (Uorc.Core.RuntimeAdmission::addOwnedRoot (project state lease) (Uorc.Core.RuntimeAdmission::addOwnedRoot (project (project state parameterTypes) root) (Uorc.Core.RuntimeAdmission::addOwnedRoot (project (project state localTypes) root) (Uorc.Core.RuntimeAdmission::addOwnedRoot (project (project state instructions) root) (Uorc.Core.RuntimeAdmission::addOwnedRoot (project (project state lastUses) root) (Uorc.Core.RuntimeAdmission::addOwnedRoot (project (project state lifetimes) root) (Uorc.Core.RuntimeAdmission::addOwnedRoot (project (project state capturesA) root) (Uorc.Core.RuntimeAdmission::addOwnedRoot (project (project state capturesB) root) (Uorc.Core.RuntimeAdmission::addOwnedRoot (project (project state bucketSlots) root) (Uorc.Core.RuntimeAdmission::addOwnedRoot (project state initialRelease) (Uorc.Core.RuntimeAdmission::addOwnedRoot (project state blockRecord) (Uorc.Core.RuntimeAdmission::addOwnedRoot (project state instructionRecord) (Uorc.Core.RuntimeAdmission::addOwnedRoot (project state bucketRecord) (Uorc.Core.RuntimeAdmission::addOwnedRoot (project state literal) (Uorc.Core.RuntimeAdmission::addOwnedRoot (project state auxFirst) (Uorc.Core.RuntimeAdmission::addOwnedRoot (project state auxSecond) (Uorc.Core.RuntimeAdmission::addOwnedRoot (project state control) (Uorc.Core.RuntimeByteBuilder::builderRoots (project state literalBuilder) (Uorc.Core.RuntimeProgram::rootInventory (project state program) caller))))))))))))))))))),
  axioms: [],
  executable: true
}

definition{
  name: countHandle,
  parameters: [
    roots : (List Uorc.Core.RuntimeValue::Handle),
    selected : Uorc.Core.RuntimeValue::Handle
  ],
  result: UInt64,
  body: (match roots
    | List.nil => (0 : UInt64)
    | List.cons head tail => (Uorc.Core.Heap::wordAdd (if (prim equal (project head id) (project selected id) : Bool) then (1 : UInt64) else (0 : UInt64)) (countHandle tail selected))),
  axioms: [Quot.sound, propext],
  executable: true,
  recursive_argument: roots
}

definition{
  name: checkRootTickets,
  parameters: [
    arena : Uorc.Core.Heap::Arena,
    allRoots : (List Uorc.Core.RuntimeValue::Handle),
    remaining : (List Uorc.Core.RuntimeValue::Handle)
  ],
  result: Bool,
  body: (match remaining
    | List.nil => true
    | List.cons head tail => (if (Uorc.Core.RuntimeValue::isNull head) then (checkRootTickets arena allRoots tail) else match{
      branches: [
        {body: false, binders: [heapFailure], constructor: Result.error},
        {body: (checkRootTickets arena allRoots tail), binders: [checked], constructor: Result.ok}
      ],
      scrutinee: (Uorc.Core.Heap::checkMoveRoot arena head (countHandle allRoots head))
    })),
  axioms: [Quot.sound, propext],
  executable: true,
  recursive_argument: remaining
}

definition{
  name: restoreInternal,
  parameters: [
    state : AwaitingState,
    arena : Uorc.Core.Heap::Arena,
    usage : Uorc.Core.RuntimeResources::Usage,
    caller : Uorc.Core.RuntimeProgram::CallerRoots
  ],
  result: Uorc.Core.RuntimeAdmission::Cursor,
  body: (record Uorc.Core.RuntimeAdmission::Cursor {arena := arena, usage := usage, budget := (project state budget), caller := caller, program := (project state program), order := (project state order), pendingOrder := (project state pendingOrder), pendingEvent := (project state pendingEvent), instruction := (project state instruction), parameterTypes := (project state parameterTypes), localTypes := (project state localTypes), instructions := (project state instructions), lastUses := (project state lastUses), lifetimes := (project state lifetimes), capturesA := (project state capturesA), capturesB := (project state capturesB), bucketSlots := (project state bucketSlots), initialRelease := (project state initialRelease), blockRecord := (project state blockRecord), instructionRecord := (project state instructionRecord), bucketRecord := (project state bucketRecord), literal := (project state literal), auxFirst := (project state auxFirst), auxSecond := (project state auxSecond), control := (project state control), declaredResult := (project state declaredResult), returnSlot := (project state returnSlot), lifetimeCursor := (project state lifetimeCursor), pendingSlot := (project state pendingSlot), bucketIndex := (project state bucketIndex), bucketCount := (project state bucketCount), bucketInitial := (project state bucketInitial), task := (project state task), failure := (project state failure), literalBuilder := (project state literalBuilder), literalOffset := (project state literalOffset), literalRemaining := (project state literalRemaining)}),
  axioms: [],
  executable: true
}

definition{
  name: take,
  parameters: [cursor : Uorc.Core.RuntimeAdmission::Cursor, role : ProducerRole],
  result: TakeEvent,
  body: (if ((((((true && match{
    branches: [
      {
        body: (match stage
          | Uorc.Core.RuntimeAdmission::Stage.initialize => false
          | Uorc.Core.RuntimeAdmission::Stage.awaiting => true
          | Uorc.Core.RuntimeAdmission::Stage.dispatch => false
          | Uorc.Core.RuntimeAdmission::Stage.ack => false
          | Uorc.Core.RuntimeAdmission::Stage.parameterType => false
          | Uorc.Core.RuntimeAdmission::Stage.snapshotParameters => false
          | Uorc.Core.RuntimeAdmission::Stage.scalar => false
          | Uorc.Core.RuntimeAdmission::Stage.capture => false
          | Uorc.Core.RuntimeAdmission::Stage.captureUse => false
          | Uorc.Core.RuntimeAdmission::Stage.literalStart => false
          | Uorc.Core.RuntimeAdmission::Stage.literalSpan => false
          | Uorc.Core.RuntimeAdmission::Stage.literalRead => false
          | Uorc.Core.RuntimeAdmission::Stage.literalFinish => false
          | Uorc.Core.RuntimeAdmission::Stage.literalJoin => false
          | Uorc.Core.RuntimeAdmission::Stage.literalJoinDone => false
          | Uorc.Core.RuntimeAdmission::Stage.literalDropLeaf => false
          | Uorc.Core.RuntimeAdmission::Stage.instructionField index => false
          | Uorc.Core.RuntimeAdmission::Stage.instructionAttach => false
          | Uorc.Core.RuntimeAdmission::Stage.instructionType => false
          | Uorc.Core.RuntimeAdmission::Stage.instructionCleanup index => false
          | Uorc.Core.RuntimeAdmission::Stage.blockReturn => false
          | Uorc.Core.RuntimeAdmission::Stage.lifetime => false
          | Uorc.Core.RuntimeAdmission::Stage.bucketAppend => false
          | Uorc.Core.RuntimeAdmission::Stage.bucketField index => false
          | Uorc.Core.RuntimeAdmission::Stage.bucketAttach => false
          | Uorc.Core.RuntimeAdmission::Stage.bucketCleanup => false
          | Uorc.Core.RuntimeAdmission::Stage.bucketCleanupSlots => false
          | Uorc.Core.RuntimeAdmission::Stage.blockField index => false
          | Uorc.Core.RuntimeAdmission::Stage.blockAttach => false
          | Uorc.Core.RuntimeAdmission::Stage.blockCount => false
          | Uorc.Core.RuntimeAdmission::Stage.blockCleanup index => false
          | Uorc.Core.RuntimeAdmission::Stage.sealCheck => false
          | Uorc.Core.RuntimeAdmission::Stage.sealFinish => false
          | Uorc.Core.RuntimeAdmission::Stage.ready => false),
        binders: [stage],
        constructor: Uorc.Core.RuntimeAdmission::Task.stage
      },
      {
        body: false,
        binders: [progress, slot, after],
        constructor: Uorc.Core.RuntimeAdmission::Task.table
      },
      {
        body: false,
        binders: [record, slot, after],
        constructor: Uorc.Core.RuntimeAdmission::Task.record
      },
      {
        body: false,
        binders: [old, after],
        constructor: Uorc.Core.RuntimeAdmission::Task.beginRelease
      },
      {body: false, binders: [after], constructor: Uorc.Core.RuntimeAdmission::Task.draining},
      {
        body: false,
        binders: [operation, after],
        constructor: Uorc.Core.RuntimeAdmission::Task.sequence
      },
      {
        body: false,
        binders: [builder, finishing, after],
        constructor: Uorc.Core.RuntimeAdmission::Task.builder
      }
    ],
    scrutinee: (project cursor task)
  }) && (match (project cursor failure)
    | Option.none => true
    | Option.some failure => false)) && (match (project cursor pendingEvent)
    | Option.none => true
    | Option.some event => false)) && (match (project cursor pendingOrder)
    | Option.none => true
    | Option.some order => false)) && (not (Uorc.Core.Heap::isReleasing (project cursor arena)))) && (not (Uorc.Core.RuntimeValue::isNull (project cursor control)))) then (match (Uorc.Core.RuntimeResources::charge (project cursor budget) (project cursor usage) (record Uorc.Core.RuntimeResources::Event {work := (Uorc.Core.RuntimeWork::fee (ctor Uorc.Core.RuntimeWork::Operation.boundaryTake)), emittedBytes := (0 : UInt64), enteredFrames := (0 : UInt64), exitedFrames := (0 : UInt64)}))
    | Result.error resourceError => (ctor TakeEvent.failed (ctor Error.resource resourceError) cursor)
    | Result.ok charged => (let roots : (List Uorc.Core.RuntimeValue::Handle) := (Uorc.Core.RuntimeAdmission::rootInventory cursor); (if (checkRootTickets (project cursor arena) roots roots) then (match (Uorc.Core.Heap::allocateControl (project cursor arena) (64 : UInt64))
      | Result.error heapFailure => (ctor TakeEvent.failed (ctor Error.heap heapFailure) (Uorc.Core.RuntimeAdmission::with_usage cursor charged))
      | Result.ok allocation => (match (Uorc.Core.RuntimeResources::observeArena (project cursor budget) charged (project allocation arena))
        | Result.error resourceError => (ctor TakeEvent.retainedFailure (ctor Error.resource resourceError) (record Taken {state := (record AwaitingState {budget := (project cursor budget), program := (project cursor program), order := (project cursor order), pendingOrder := (project cursor pendingOrder), pendingEvent := (project cursor pendingEvent), instruction := (project cursor instruction), parameterTypes := (project cursor parameterTypes), localTypes := (project cursor localTypes), instructions := (project cursor instructions), lastUses := (project cursor lastUses), lifetimes := (project cursor lifetimes), capturesA := (project cursor capturesA), capturesB := (project cursor capturesB), bucketSlots := (project cursor bucketSlots), initialRelease := (project cursor initialRelease), blockRecord := (project cursor blockRecord), instructionRecord := (project cursor instructionRecord), bucketRecord := (project cursor bucketRecord), literal := (project cursor literal), auxFirst := (project cursor auxFirst), auxSecond := (project cursor auxSecond), control := (project cursor control), declaredResult := (project cursor declaredResult), returnSlot := (project cursor returnSlot), lifetimeCursor := (project cursor lifetimeCursor), pendingSlot := (project cursor pendingSlot), bucketIndex := (project cursor bucketIndex), bucketCount := (project cursor bucketCount), bucketInitial := (project cursor bucketInitial), task := (project cursor task), failure := (project cursor failure), literalBuilder := (project cursor literalBuilder), literalOffset := (project cursor literalOffset), literalRemaining := (project cursor literalRemaining), role := role, lease := (project allocation handle), expectedCaller := (project cursor caller), minimumNextHandle := (project (project allocation arena) nextHandle), wasExhausted := (project (project allocation arena) exhausted), minimumWork := (project charged work), minimumOutput := (project charged outputBytes), expectedFrames := (project charged frames), minimumPeak := (project charged peakOwnedBytes)}), arena := (project allocation arena), usage := charged, caller := (project cursor caller)}))
        | Result.ok observed => (ctor TakeEvent.taken (record Taken {state := (record AwaitingState {budget := (project cursor budget), program := (project cursor program), order := (project cursor order), pendingOrder := (project cursor pendingOrder), pendingEvent := (project cursor pendingEvent), instruction := (project cursor instruction), parameterTypes := (project cursor parameterTypes), localTypes := (project cursor localTypes), instructions := (project cursor instructions), lastUses := (project cursor lastUses), lifetimes := (project cursor lifetimes), capturesA := (project cursor capturesA), capturesB := (project cursor capturesB), bucketSlots := (project cursor bucketSlots), initialRelease := (project cursor initialRelease), blockRecord := (project cursor blockRecord), instructionRecord := (project cursor instructionRecord), bucketRecord := (project cursor bucketRecord), literal := (project cursor literal), auxFirst := (project cursor auxFirst), auxSecond := (project cursor auxSecond), control := (project cursor control), declaredResult := (project cursor declaredResult), returnSlot := (project cursor returnSlot), lifetimeCursor := (project cursor lifetimeCursor), pendingSlot := (project cursor pendingSlot), bucketIndex := (project cursor bucketIndex), bucketCount := (project cursor bucketCount), bucketInitial := (project cursor bucketInitial), task := (project cursor task), failure := (project cursor failure), literalBuilder := (project cursor literalBuilder), literalOffset := (project cursor literalOffset), literalRemaining := (project cursor literalRemaining), role := role, lease := (project allocation handle), expectedCaller := (project cursor caller), minimumNextHandle := (project (project allocation arena) nextHandle), wasExhausted := (project (project allocation arena) exhausted), minimumWork := (project observed work), minimumOutput := (project observed outputBytes), expectedFrames := (project observed frames), minimumPeak := (project observed peakOwnedBytes)}), arena := (project allocation arena), usage := observed, caller := (project cursor caller)})))) else (ctor TakeEvent.failed (ctor Error.invalidTicketInventory) (Uorc.Core.RuntimeAdmission::with_usage cursor charged))))) else (ctor TakeEvent.failed (ctor Error.busy) cursor)),
  axioms: [Quot.sound, propext],
  executable: true
}

definition{
  name: leaseValid,
  parameters: [state : AwaitingState, arena : Uorc.Core.Heap::Arena],
  result: Bool,
  body: (match (Uorc.Core.Heap::lookup arena (project state lease))
    | Result.error heapFailure => false
    | Result.ok entry => (((true && (prim equal (project entry roots) (1 : UInt64) : Bool)) && (prim equal (project entry incoming) (1 : UInt64) : Bool)) && (match (project entry node)
      | Uorc.Core.Heap::Node.image field0 field1 => false
      | Uorc.Core.Heap::Node.sequenceLeaf field0 field1 field2 => false
      | Uorc.Core.Heap::Node.sequenceBranch field0 field1 field2 field3 => false
      | Uorc.Core.Heap::Node.storeLeaf field0 => false
      | Uorc.Core.Heap::Node.storeBranch field0 field1 => false
      | Uorc.Core.Heap::Node.slot field0 field1 => false
      | Uorc.Core.Heap::Node.rootBranch field0 field1 field2 => false
      | Uorc.Core.Heap::Node.metadataWord field0 field1 field2 => false
      | Uorc.Core.Heap::Node.metadataLink field0 field1 field2 => false
      | Uorc.Core.Heap::Node.metadataBranch field0 field1 field2 => false
      | Uorc.Core.Heap::Node.ownerRoot field0 field1 => false
      | Uorc.Core.Heap::Node.imageReservation field0 field1 => ((true && (match field0
        | Uorc.Core.Heap::ImageKind.input => false
        | Uorc.Core.Heap::ImageKind.archive => false
        | Uorc.Core.Heap::ImageKind.basis => false
        | Uorc.Core.Heap::ImageKind.runtime => false
        | Uorc.Core.Heap::ImageKind.control => true)) && (prim equal field1 (64 : UInt64) : Bool))))),
  axioms: [Quot.sound, propext],
  executable: true
}

definition{
  name: reattach,
  parameters: [
    state : AwaitingState,
    role : ProducerRole,
    arena : Uorc.Core.Heap::Arena,
    usage : Uorc.Core.RuntimeResources::Usage,
    caller : Uorc.Core.RuntimeProgram::CallerRoots
  ],
  result: ReturnEvent,
  body: (let taken : Taken := (record Taken {state := state, arena := arena, usage := usage, caller := caller}); (if (not (sameRole (project state role) role)) then (ctor ReturnEvent.failed (ctor Error.wrongRole) (record ReturnCursor {taken := taken, leaseLive := true})) else (if (Uorc.Core.Heap::isReleasing arena) then (ctor ReturnEvent.failed (ctor Error.foreignRelease) (record ReturnCursor {taken := taken, leaseLive := true})) else (if (not (sameNonProducerCaller role (project state expectedCaller) caller)) then (ctor ReturnEvent.failed (ctor Error.callerMismatch) (record ReturnCursor {taken := taken, leaseLive := true})) else (if (not (((((((true && (not (Uorc.Core.Heap::less (project arena nextHandle) (project state minimumNextHandle)))) && (if (project state wasExhausted) then (project arena exhausted) else true)) && (not (Uorc.Core.Heap::less (project usage work) (project state minimumWork)))) && (not (Uorc.Core.Heap::less (project usage outputBytes) (project state minimumOutput)))) && (prim equal (project usage frames) (project state expectedFrames) : Bool)) && (not (Uorc.Core.Heap::less (project usage peakOwnedBytes) (project state minimumPeak)))) && (prim equal (project usage ownedBytes) (project arena ownedBytes) : Bool))) then (ctor ReturnEvent.failed (ctor Error.counterRegression) (record ReturnCursor {taken := taken, leaseLive := true})) else (if (not (leaseValid state arena)) then (ctor ReturnEvent.failed (ctor Error.staleSession) (record ReturnCursor {taken := taken, leaseLive := true})) else (match (Uorc.Core.RuntimeResources::charge (project state budget) usage (record Uorc.Core.RuntimeResources::Event {work := (Uorc.Core.RuntimeWork::fee (ctor Uorc.Core.RuntimeWork::Operation.boundaryReattach)), emittedBytes := (0 : UInt64), enteredFrames := (0 : UInt64), exitedFrames := (0 : UInt64)}))
    | Result.error resourceError => (ctor ReturnEvent.failed (ctor Error.resource resourceError) (record ReturnCursor {taken := taken, leaseLive := true}))
    | Result.ok charged => (let roots : (List Uorc.Core.RuntimeValue::Handle) := (rootInventory state caller); (if (checkRootTickets arena roots roots) then (match (Uorc.Core.Heap::release arena (project state lease))
      | Result.error heapFailure => (ctor ReturnEvent.failed (ctor Error.heap heapFailure) (record ReturnCursor {taken := (record Taken {state := state, arena := arena, usage := charged, caller := caller}), leaseLive := true}))
      | Result.ok transition => (ctor ReturnEvent.pending (record ReturnCursor {taken := (record Taken {state := state, arena := (project transition arena), usage := charged, caller := caller}), leaseLive := false}))) else (ctor ReturnEvent.failed (ctor Error.invalidTicketInventory) (record ReturnCursor {taken := (record Taken {state := state, arena := arena, usage := charged, caller := caller}), leaseLive := true}))))))))))),
  axioms: [Quot.sound, propext],
  executable: true
}

definition{
  name: rootInventoryWithoutLease,
  parameters: [state : AwaitingState, caller : Uorc.Core.RuntimeProgram::CallerRoots],
  result: (List Uorc.Core.RuntimeValue::Handle),
  body: (Uorc.Core.RuntimeAdmission::addOwnedRoot (project (project state parameterTypes) root) (Uorc.Core.RuntimeAdmission::addOwnedRoot (project (project state localTypes) root) (Uorc.Core.RuntimeAdmission::addOwnedRoot (project (project state instructions) root) (Uorc.Core.RuntimeAdmission::addOwnedRoot (project (project state lastUses) root) (Uorc.Core.RuntimeAdmission::addOwnedRoot (project (project state lifetimes) root) (Uorc.Core.RuntimeAdmission::addOwnedRoot (project (project state capturesA) root) (Uorc.Core.RuntimeAdmission::addOwnedRoot (project (project state capturesB) root) (Uorc.Core.RuntimeAdmission::addOwnedRoot (project (project state bucketSlots) root) (Uorc.Core.RuntimeAdmission::addOwnedRoot (project state initialRelease) (Uorc.Core.RuntimeAdmission::addOwnedRoot (project state blockRecord) (Uorc.Core.RuntimeAdmission::addOwnedRoot (project state instructionRecord) (Uorc.Core.RuntimeAdmission::addOwnedRoot (project state bucketRecord) (Uorc.Core.RuntimeAdmission::addOwnedRoot (project state literal) (Uorc.Core.RuntimeAdmission::addOwnedRoot (project state auxFirst) (Uorc.Core.RuntimeAdmission::addOwnedRoot (project state auxSecond) (Uorc.Core.RuntimeAdmission::addOwnedRoot (project state control) (Uorc.Core.RuntimeByteBuilder::builderRoots (project state literalBuilder) (Uorc.Core.RuntimeProgram::rootInventory (project state program) caller)))))))))))))))))),
  axioms: [],
  executable: true
}

definition{
  name: returnStep,
  parameters: [returning : ReturnCursor],
  result: ReturnEvent,
  body: (if (project returning leaseLive) then (ctor ReturnEvent.failed (ctor Error.busy) returning) else (match (Uorc.Core.RuntimeResources::charge (project (project (project returning taken) state) budget) (project (project returning taken) usage) (record Uorc.Core.RuntimeResources::Event {work := (Uorc.Core.RuntimeWork::fee (ctor Uorc.Core.RuntimeWork::Operation.boundaryReturn)), emittedBytes := (0 : UInt64), enteredFrames := (0 : UInt64), exitedFrames := (0 : UInt64)}))
    | Result.error resourceError => (ctor ReturnEvent.failed (ctor Error.resource resourceError) returning)
    | Result.ok charged => (match (Uorc.Core.Heap::releaseStep (project (project returning taken) arena))
      | Result.error heapFailure => (ctor ReturnEvent.failed (ctor Error.heap heapFailure) (record ReturnCursor {taken := (record Taken {state := (project (project returning taken) state), arena := (project (project returning taken) arena), usage := charged, caller := (project (project returning taken) caller)}), leaseLive := false}))
      | Result.ok transition => (if (project transition complete) then (match (Uorc.Core.RuntimeResources::observeArena (project (project (project returning taken) state) budget) charged (project transition arena))
        | Result.error resourceError => (ctor ReturnEvent.failed (ctor Error.resource resourceError) (record ReturnCursor {taken := (record Taken {state := (project (project returning taken) state), arena := (project transition arena), usage := charged, caller := (project (project returning taken) caller)}), leaseLive := false}))
        | Result.ok observed => (ctor ReturnEvent.restored (restoreInternal (project (project returning taken) state) (project transition arena) observed (project (project returning taken) caller)))) else (ctor ReturnEvent.failed (ctor Error.staleSession) (record ReturnCursor {taken := (record Taken {state := (project (project returning taken) state), arena := (project transition arena), usage := charged, caller := (project (project returning taken) caller)}), leaseLive := false})))))),
  axioms: [Quot.sound, propext],
  executable: true
}

definition{
  name: returnRootInventory,
  parameters: [returning : ReturnCursor],
  result: (List Uorc.Core.RuntimeValue::Handle),
  body: (if (project returning leaseLive) then (rootInventory (project (project returning taken) state) (project (project returning taken) caller)) else (rootInventoryWithoutLease (project (project returning taken) state) (project (project returning taken) caller))),
  axioms: [],
  executable: true
}

definition{
  name: takeWork,
  parameters: [],
  result: UInt64,
  body: (Uorc.Core.RuntimeWork::fee (ctor Uorc.Core.RuntimeWork::Operation.boundaryTake)),
  axioms: [],
  executable: true
}

definition{
  name: reattachWork,
  parameters: [],
  result: UInt64,
  body: (Uorc.Core.RuntimeWork::fee (ctor Uorc.Core.RuntimeWork::Operation.boundaryReattach)),
  axioms: [],
  executable: true
}

definition{
  name: returnStepWork,
  parameters: [],
  result: UInt64,
  body: (Uorc.Core.RuntimeWork::fee (ctor Uorc.Core.RuntimeWork::Operation.boundaryReturn)),
  axioms: [],
  executable: true
}
\end{semanticmodule}
\end{lexlean}
